\section{数值模拟与时间反演的直觉}
\label{sec:06-Random-Processes/06-Diffusion-and-SDE/04}

一个训练好的扩散模型，采样步数从 1000 降到 50，速度快了二十倍，出来的图明显糊了。可要是你只想统计生成图片的平均亮度，50 步给的数字和 1000 步几乎没有差别。

同一个离散化，同一份权重，一个指标崩了，另一个纹丝不动。这不是玄学。数值格式对``单条路径准不准''和``分布准不准''给的是两套完全不同的保证，混着谈就会得出互相矛盾的结论。

上一节的结论是解析解基本指望不上，所以这一节把方程交给计算机：先写出最朴素的离散格式，再把上面那两套保证分清楚，最后朝一个反直觉的方向走——把整个方程倒过来跑。扩散生成模型的数学入口就在那里。

\subsection{把增量公式抄成递推}

面对 \(dX_t=\mu(X_t)\,dt+\sigma(X_t)\,dW_t\)，取时间网格 \(t_k=k\Delta t\)，把微分记号当成极短时间内的真实增量直接抄下来：

\[
X_{k+1}=X_k+\mu(X_k)\,\Delta t+\sigma(X_k)\sqrt{\Delta t}\,Z_k,
\qquad Z_k\sim\mathcal N(0,1)\ \text{独立同分布}.
\]

这是 Euler--Maruyama 格式。系数在区间左端点取值，与 Itô 约定一致；换成右端点或中点，模拟的就是另一个方程了。

噪声项前面是 \(\sqrt{\Delta t}\) 而不是 \(\Delta t\)，这是本单元第一节那条标度关系的回访。那里是从离散走向连续，让 \(\Delta x=\sqrt{\Delta t}\) 保住极限；这里是从连续退回离散，同一个平方根保证离散过程的方差累积对得上。正向是收敛，反向是逼近，碰头的地方就是这个根号。程序里最常见的一类错误正是把它写成 \(\Delta t\)——跑出来的路径看着还挺像回事，方差却小了整整一个量级。

\subsection{路径错了，分布可以是对的}

误差这个词在随机的世界里裂成了两半。

强收敛盯着单条路径。固定同一份噪声实现，比较数值路径与真实路径的偏离，看 \(\E\abs{X_T-X_T^{(\Delta t)}}\) 随 \(\Delta t\) 下降得多快。系数足够光滑等常规条件下，Euler--Maruyama 的强阶是 \(1/2\)。

弱收敛只盯着分布。不管路径是否重合，只看统计量对不对，即 \(\abs{\E[f(X_T)]-\E[f(X_T^{(\Delta t)})]}\) 的下降速度。同样条件下弱阶是 \(1\)。

\begin{center}
\begin{tikzpicture}[x=1cm,y=1cm,font=\small,>=stealth]
  \draw[black!45] (0,-0.1) -- (3.7,-0.1);
  \draw[black!80,thick] (0,0.9) -- (0.3,1.15) -- (0.6,0.95) -- (0.9,1.35) -- (1.2,1.15)
    -- (1.5,1.5) -- (1.8,1.25) -- (2.1,1.6) -- (2.4,1.35) -- (2.7,1.75) -- (3.0,1.5)
    -- (3.3,1.85) -- (3.6,1.65);
  \draw[blue!70,thick,dashed] (0,0.9) -- (0.6,1.25) -- (1.2,0.95) -- (1.8,1.45) -- (2.4,1.0)
    -- (3.0,1.35) -- (3.6,0.75);
  \node[right,font=\footnotesize,black!80] at (0.15,2.35) {真实路径};
  \node[right,font=\footnotesize,blue!70] at (0.15,2.0) {大步长的数值路径};
  \node[below,font=\footnotesize,align=center] at (1.8,-0.25) {同一份噪声下\\两条路径分道扬镳};
  \draw[black!25] (4.4,-0.6) -- (4.4,2.6);
  \begin{scope}[shift={(5.1,0)}]
    \draw[black!45] (-0.1,-0.1) -- (3.9,-0.1);
    \foreach \i/\h in {0/0.10,1/0.22,2/0.46,3/0.86,4/1.30,5/1.68,6/1.80,7/1.52,8/1.06,
                       9/0.62,10/0.30,11/0.12}
      \fill[black!15] ({0.3*\i},-0.1) rectangle ({0.3*\i+0.3},{\h-0.1});
    \draw[blue!75,thick] (0,0.02) -- (0.3,0.02) -- (0.3,0.16) -- (0.6,0.16) -- (0.6,0.36)
      -- (0.9,0.36) -- (0.9,0.82) -- (1.2,0.82) -- (1.2,1.24) -- (1.5,1.24) -- (1.5,1.62)
      -- (1.8,1.62) -- (1.8,1.74) -- (2.1,1.74) -- (2.1,1.46) -- (2.4,1.46) -- (2.4,1.12)
      -- (2.7,1.12) -- (2.7,0.56) -- (3.0,0.56) -- (3.0,0.26) -- (3.3,0.26) -- (3.3,0.08)
      -- (3.6,0.08);
    \node[right,font=\footnotesize,black!55] at (0.05,2.35) {真实分布};
    \node[right,font=\footnotesize,blue!75] at (0.05,2.0) {大步长的数值分布};
    \node[below,font=\footnotesize,align=center] at (1.8,-0.25) {同一个步长下\\两个直方图几乎重合};
  \end{scope}
\end{tikzpicture}

\emph{左右两张图来自同一次实验、同一个步长。左边逐条比对，数值路径早早偏离，误差按 \(\Delta t^{1/2}\) 收窄，慢得让人不安；右边把上万条路径的终点堆成直方图，两条分布几乎压在一起，误差按 \(\Delta t\) 收窄。问``模拟准不准''之前得先说清楚要拿它做什么。}
\end{center}

强阶掉到 \(1/2\)，原因还是那个平方根。确定性的 Euler 格式每步丢掉的是 \(O(\Delta t^2)\) 的二阶项，累积成 \(O(\Delta t)\)。随机情形下，把 \(\sigma\) 在整个小区间上冻结为 \(\sigma(X_k)\) 这个动作，漏掉的是形如 \(\sigma\sigma'\int(W_s-W_{t_k})\,dW_s\) 的一项，量级 \(\Delta t\) 而不是 \(\Delta t^2\)——两个 \(\sqrt{\Delta t}\) 相乘的结果。Milstein 格式把这一项补回去，加上 \(\tfrac12\sigma\sigma'\big((\Delta W_k)^2-\Delta t\big)\) 就把强阶提到 \(1\)。补的正是二次变差贡献的那部分，而当 \(\sigma\) 与状态无关时 \(\sigma'=0\)，这项自动消失，Euler--Maruyama 本来就是强一阶的。扩散生成模型用的噪声调度恰好属于这一类。

误差还有第二个来源。跑一次模拟只得到一条路径，``看着像''不能当证据。估计 \(\E[f(X_T)]\) 要独立跑 \(M\) 条取平均，大数定律保证收敛，中心极限定理给出涨落量级 \(\sigma_f/\sqrt M\)，正是\hyperref[ch:044]{``极限定理：重复随机性为何会产生稳定规律''}那两件工具的现场组合。于是总误差被劈成两块：离散化偏差随 \(\Delta t\) 下降，Monte Carlo 涨落随 \(M\) 下降，而计算预算是固定的。步长减半让每条路径贵一倍，样本翻倍让方差减半——怎么分这笔钱，是随机数值分析最基本的一场博弈。

\subsection{先把数据洗成噪声}

现在换个立场：不再是模拟一个已知系统，而是想造出新样本。

设数据来自某个未知分布 \(p_0\)，比如所有``自然图片''构成的那个分布。它复杂到写不出密度函数，但你能从中采样——训练集就是样本。

取一个与数据毫无关系的固定加噪过程，比如上一节反复出现的均值回复模型 \(dX_t=-\theta X_t\,dt+\sigma\,dW_t\)，让它从数据样本出发跑起来。前面已经算过：初值按 \(e^{-\theta t}\) 被遗忘，任何起点的分布都被拖向同一个 Gaussian。所以不论 \(X_0\) 落在 \(p_0\) 的哪个角落，跑够久之后它的分布都近似于同一个已知的、极易采样的噪声分布。

用密度视角重述一遍（见\hyperref[sec:06-Random-Processes/06-Diffusion-and-SDE/02]{``Fokker--Planck 方程描述密度的流动''}）：漂移把质量往中心搬，扩散把质量摊平，数据密度里那些精细结构一层层被抹掉，终态是纯噪声。这一步机械得可以写在纸上，不含任何待学的参数。

问题在于它是单向的。信息在往外漏，怎么可能倒回去？

\subsection{倒着跑，要的是得分}

在常规正则条件下有一个不平凡的事实：把扩散过程的时间倒过来，得到的仍然是一个扩散过程。它有自己的 SDE，只是漂移项里多了一样东西。

考虑生成模型里最常用的情形——扩散强度 \(\sigma(t)\) 可随时间变化但不依赖状态，且各方向相同。若正向方程是 \(dX=f(X,t)\,dt+\sigma(t)\,dW\)，则反向方程形如

\[
dX=\Big(f(X,t)-\sigma(t)^2\,\nabla_x\log p_t(X)\Big)dt+\sigma(t)\,d\bar W,
\]

时间从 \(T\) 流向 \(0\)，\(\bar W\) 是反向时间里的另一个布朗运动。多出来的 \(\nabla_x\log p_t(x)\) 叫得分（score），是密度对数的梯度，逐点指向密度上升最快的方向。

\begin{center}
\begin{tikzpicture}[x=1cm,y=1cm,font=\small,>=stealth]
  \draw[->,black!70,thick] (0.1,1.95) -- (10.2,1.95);
  \node[above,font=\footnotesize,black!75] at (5.15,2.0) {前向 SDE：加噪，机械执行，无参数可学};
  \draw[->,red!70,thick] (10.2,-1.25) -- (0.1,-1.25);
  \node[below,font=\footnotesize,red!75] at (5.15,-1.32)
    {反向 SDE：去噪，漂移里唯一的未知量是得分 \(\nabla_x\log p_t\)};
  \begin{scope}
    \draw[black!30] (0,-0.15) rectangle (2.2,1.55);
    \draw[blue!75,thick] plot[domain=0.05:2.15,samples=70]
      (\x,{-0.1+0.95*exp(-(\x-0.6)*(\x-0.6)/0.025)+0.8*exp(-(\x-1.6)*(\x-1.6)/0.03)});
    \node[below,font=\footnotesize] at (1.1,-0.25) {\(p_0\)：数据};
  \end{scope}
  \begin{scope}[shift={(2.7,0)}]
    \draw[black!30] (0,-0.15) rectangle (2.2,1.55);
    \draw[blue!75,thick] plot[domain=0.05:2.15,samples=70]
      (\x,{-0.1+0.72*exp(-(\x-0.65)*(\x-0.65)/0.12)+0.62*exp(-(\x-1.58)*(\x-1.58)/0.13)});
    \node[below,font=\footnotesize] at (1.1,-0.25) {\(p_{t_1}\)};
  \end{scope}
  \begin{scope}[shift={(5.4,0)}]
    \draw[black!30] (0,-0.15) rectangle (2.2,1.55);
    \draw[blue!75,thick] plot[domain=0.05:2.15,samples=70]
      (\x,{-0.1+0.58*exp(-(\x-0.9)*(\x-0.9)/0.40)+0.40*exp(-(\x-1.6)*(\x-1.6)/0.38)});
    \node[below,font=\footnotesize] at (1.1,-0.25) {\(p_{t_2}\)};
  \end{scope}
  \begin{scope}[shift={(8.1,0)}]
    \draw[black!30] (0,-0.15) rectangle (2.2,1.55);
    \draw[blue!75,thick] plot[domain=0.05:2.15,samples=70]
      (\x,{-0.1+0.68*exp(-(\x-1.1)*(\x-1.1)/0.70)});
    \node[below,font=\footnotesize] at (1.1,-0.25) {\(p_T\)：纯噪声};
  \end{scope}
\end{tikzpicture}

\emph{扩散生成模型的全景。上排箭头是前向过程，数据密度的多峰结构被逐帧抹平，最后塌成一个已知的 Gaussian；这一路不需要训练。下排箭头把同一条时间轴倒过来走，从纯噪声出发重建数据密度，而整个反向方程里唯一不能直接写出的东西就是得分。学会了得分，就学会了生成——这张图是本单元与现代生成模型之间的全部接口。}
\end{center}

多出这一项的道理并不神秘。纯扩散只会把概率摊平，粒子从密集处走向稀疏处。想让已经摊开的噪声重新收拢成数据的形状，就得额外施加一股把粒子往高密度区域推的力，而``哪边密度更高''这个信息，正好就是 \(\nabla_x\log p_t\)。系数 \(\sigma(t)^2\) 说明这股推力要和当时的扩散强度成正比：噪声撒得越猛，往回拽得越用力。

得分本身随时间变化，这一点决定了生成过程的观感。\(t\) 大、噪声重时，\(p_t\) 已被抹得很平滑，得分场只能给出粗线条的指向，说的是``大致该往哪个区域去''；\(t\) 小、噪声轻时，\(p_t\) 恢复了精细结构，得分场开始分辨边界与纹理。反向过程因此是先勾轮廓再补细节。这和上一节退火采样那张双阱图讲的是同一件事：高温先找区域，降温再压细节，只不过这里的``温度''换成了噪声水平，而且每个水平上的漂移都是学出来的。

若扩散矩阵依赖状态，反向漂移一般还要多一个扩散矩阵的散度项，不能照搬上式。但主干不变：反演的门票就是得分。

\subsection{三步拼成一台机器}

扩散生成模型把上面两段拼起来。第一步前向加噪，用固定 SDE 把数据分布洗成近似标准 Gaussian，纯机械，无参可学。第二步学得分，用数据估计各噪声水平上的得分场 \(\nabla_x\log p_t\)，这是一个回归问题，交给神经网络，全程不需要写出 \(p_0\) 的密度。第三步反向生成，从纯噪声采样，把估计出的得分代进反向 SDE 的漂移，沿反向时间用 Euler--Maruyama 一步步走，终点样本近似服从 \(p_0\)。

聪明的地方在于绕开了什么。对图片这类分布，写出密度基本无望，配分函数更是算不动；而反向 SDE 从头到尾只需要密度的对数梯度。这正是上一节末尾 Langevin 采样绕开归一化常数的同一个理由，只是用在了一个动态版本上。上一节的存在唯一性条件保证正向方程写得合法，本节的时间反演保证倒着跑的方向是对的，剩下的全部工程任务就是把得分估准——这件事留给下一节。

还有一个值得知道的变体。反向 SDE 有一个共享同样时刻边缘分布的确定性对应物，把随机项去掉、把得分项的系数减半，得到的常微分方程

\[
\frac{dX}{dt}=f(X,t)-\frac12\sigma(t)^2\nabla_x\log p_t(X)
\]

沿时间倒推，每个时刻的分布与随机版本完全一致。它让采样变得可复现，允许用成熟的 ODE 求解器大步走，实践中的快速采样器多半基于它。同一族边缘分布可以由不止一个过程生成——这件事本身就值得留意：分布层面的等价并不意味着路径层面的等价。

\subsection{三笔近似的账}

严格的时间反演定理需要密度与得分的正则条件，证明远超本书范围。这一节交付的是机制，不是逐条前提的定理重述。

实际生成质量的损失恰好对应三处近似。其一是得分估计误差，网络没能完全学到真实的 \(\nabla_x\log p_t\)，在样本稀疏的低密度区尤其不准。其二是反向离散的步长误差，Euler--Maruyama 每步都在做线性近似，采样步数从 1000 砍到 50 时崩掉的就是这一项——而它属于强收敛的辖区，因为你在意的是这一条路径最后长什么样，不是一万张图的平均亮度。开头那个现象到这里就解释完了。其三是终点失配，前向只跑有限时间，到达的分布与标准正态之间还差一点，反向的起点因此本就站偏了。三项各有活跃的研究在压缩，也各有各的调参手段。

从更远处看，这与\hyperref[sec:13-Machine-Learning/08-Sampling-and-Inference/03]{``变分推断把后验近似变成优化''}是同一个套路：精确对象算不出来，就换一个可优化的近似对象顶上，再让采样或推断走完最后一步。变分推断把积分换成优化，扩散生成把采样换成沿一个学来的向量场走路。两者的可信度都取决于同一件事——你是否清楚自己替掉了什么。

下一节回答那个悬着的问题：密度都不知道，得分怎么可能学得到。
